\documentclass[11pt]{article}
% ===========================================================================
%  Shared preamble for the AP Calculus AB exam-preparation worksheets.
%  Input right after \documentclass. Files beginning with "_" are not compiled
%  on their own. Navy and gold, to match the course's lesson pages.
% ===========================================================================
\usepackage[letterpaper,top=2.3cm,bottom=2.2cm,left=1.7cm,right=1.7cm,headheight=28pt,footskip=30pt]{geometry}
\usepackage{amsmath,amssymb}
\usepackage[T1]{fontenc}
\usepackage[utf8]{inputenc}
\usepackage{lmodern}
\usepackage{xcolor}
\usepackage[most]{tcolorbox}
\usepackage{enumitem}
\usepackage{fancyhdr}
\usepackage{array}
\usepackage{tikz}
\usepackage{pgfplots}
\pgfplotsset{compat=1.18}

\definecolor{navy}{HTML}{161E54}
\definecolor{ink}{HTML}{27357E}
\definecolor{gold}{HTML}{B98A00}
\definecolor{goldbg}{HTML}{FFFDF3}
\definecolor{inkbg}{HTML}{F6F7FD}
\definecolor{rule}{HTML}{C9CFF0}
\definecolor{okgreen}{HTML}{1B7A44}
\definecolor{okbg}{HTML}{F1F8F3}

\pagestyle{fancy}
\fancyhf{}
\fancyhead[L]{\footnotesize Name: \underline{\hspace{4cm}}}
\fancyhead[C]{\textbf{\color{navy}Integration Academy}\\[-2pt]{\scriptsize Calculus AB \textperiodcentered{} AP\textsuperscript{\textregistered} Exam Preparation}}
\fancyfoot[L]{\tiny\color{black!55}\parbox[t]{0.78\textwidth}{AP\textsuperscript{\textregistered} is a registered trademark of the College Board, which was not involved in the production of, and does not endorse, this material. \textcopyright{} Integration Academy.}}
\fancyfoot[R]{\scriptsize\thepage}
\renewcommand{\headrulewidth}{1.2pt}
\renewcommand{\headrule}{\hbox to\headwidth{\color{gold}\leaders\hrule height \headrulewidth\hfill}}
\renewcommand{\footrulewidth}{0pt}
\setlength{\parindent}{0pt}
\setlength{\parskip}{4pt}

% Title block: \wstitle{code}{title}{unit line}
\newcommand{\wstitle}[3]{\begin{center}{\small\bfseries\color{gold}WORKSHEET #1}\\[3pt]{\LARGE\bfseries\color{navy}#2}\\[3pt]{\small\color{black!60}#3}\end{center}\vspace{2pt}}
% Section heading
\newcommand{\sech}[1]{\par\vspace{10pt}{\large\bfseries\color{navy}#1}\par\nobreak\vspace{2pt}{\color{rule}\hrule height0.8pt}\vspace{7pt}}
% A level tag printed beside a group of questions
\newcommand{\level}[1]{\par\vspace{6pt}{\footnotesize\bfseries\color{gold}\MakeUppercase{#1}}\par\nobreak\vspace{2pt}}

\newtcolorbox{ideabox}{colframe=ink,colback=inkbg,coltitle=white,colbacktitle=ink,fonttitle=\bfseries\small,title={Key ideas},arc=2pt,boxrule=0.8pt,left=8pt,right=8pt}
\newtcolorbox{example}[1]{colframe=rule,colback=white,coltitle=navy,colbacktitle=inkbg,fonttitle=\bfseries\small,title={#1},arc=2pt,boxrule=0.8pt,borderline west={3pt}{0pt}{ink},left=10pt,right=8pt}
\newtcolorbox{frbox}[1]{colframe=gold,colback=goldbg,coltitle=white,colbacktitle=gold,fonttitle=\bfseries\small,title={#1},arc=2pt,boxrule=0.8pt,breakable,left=8pt,right=8pt}
\newtcolorbox{solbox}{colframe=okgreen,colback=okbg,coltitle=white,colbacktitle=okgreen,fonttitle=\bfseries\small,title={Worked solutions},arc=2pt,boxrule=0.8pt,breakable,left=8pt,right=8pt}
\newtcolorbox{scorebox}[1]{colframe=navy,colback=white,coltitle=white,colbacktitle=navy,fonttitle=\bfseries\small,title={#1},arc=2pt,boxrule=0.8pt,breakable,left=8pt,right=8pt}

\newcommand{\soln}{\par\smallskip\textbf{\color{ink}Solution.}\ }
\newcommand{\result}[1]{\par\smallskip\textbf{#1}}
% blank working space of a given height
\newcommand{\work}[1]{\par\vspace{#1}}
% a point value in the margin of a scoring guide
\newcommand{\pts}[1]{\hfill{\footnotesize\bfseries\color{gold}[#1]}}

% Coordinate grid for graph questions: \begin{apgraph}{xmin}{xmax}{ymin}{ymax} ... \end{apgraph}
\newenvironment{apgraph}[4]{\begin{center}\begin{tikzpicture}\begin{axis}[
  axis lines=middle, xmin=#1, xmax=#2, ymin=#3, ymax=#4,
  grid=both, grid style={black!12}, xtick distance=1, ytick distance=1,
  ticklabel style={font=\scriptsize}, xlabel={$x$}, ylabel={$y$},
  width=10cm, height=7.2cm, samples=80, clip=false,
  every axis plot/.append style={line width=1.1pt, color=ink}]}{\end{axis}\end{tikzpicture}\end{center}}
\newcommand{\closeddot}[2]{\addplot[only marks, mark=*, mark size=2.2pt, color=ink] coordinates {(#1,#2)};}
\newcommand{\opendot}[2]{\addplot[only marks, mark=*, mark size=2.2pt, color=ink, mark options={fill=white}] coordinates {(#1,#2)};}

% Multiple-choice options: four across, or stacked when the options are long.
\newcommand{\opts}[4]{\par\vspace{3pt}\begin{tabular}{@{}p{0.235\linewidth}p{0.235\linewidth}p{0.235\linewidth}p{0.235\linewidth}@{}}(A)\ #1 & (B)\ #2 & (C)\ #3 & (D)\ #4\end{tabular}\par\vspace{9pt}}
\newcommand{\optsv}[4]{\par\vspace{3pt}\begin{tabular}{@{}l@{\ \ }p{0.9\linewidth}@{}}(A) & #1\\(B) & #2\\(C) & #3\\(D) & #4\end{tabular}\par\vspace{9pt}}

\fancyhead[R]{\footnotesize\color{black!60}Worksheet 7.4}
\begin{document}

\wstitle{7.4}{Reasoning from a Slope Field}{Unit 7: Differential Equations}

\begin{ideabox}
\begin{itemize}[leftmargin=*,itemsep=2pt,topsep=2pt]
  \item An \textbf{equilibrium solution} is a constant solution: a horizontal line in the slope field.
  \item A solution curve cannot cross an equilibrium. The sign of $\dfrac{dy}{dt}$ on each side tells you whether solutions approach it or move away.
  \item $\dfrac{d^2y}{dx^2}$ is found by differentiating the equation and substituting for $\dfrac{dy}{dx}$. Its sign gives concavity.
  \item For a logistic-type equation, growth is fastest halfway between the equilibria, and the solution has a point of inflection there.
\end{itemize}
\end{ideabox}

\sech{Worked examples}
\begin{example}{Example 1 \textperiodcentered{} One equilibrium}
Describe the long-run behaviour of solutions of $\dfrac{dy}{dt}=3-y$.\soln Equilibrium $y=3$. Above it the derivative is negative; below it, positive.\result{Every solution approaches 3.}
\end{example}

\begin{example}{Example 2 \textperiodcentered{} Two equilibria}
For $\dfrac{dy}{dt}=y(5-y)$, describe the solutions with $y(0)=1$ and with $y(0)=7$.\soln Equilibria 0 and 5. Between them the derivative is positive; above 5 it is negative.\result{The first increases toward 5; the second decreases toward 5.}
\end{example}

\begin{example}{Example 3 \textperiodcentered{} A second derivative}
For $\dfrac{dy}{dx}=3x-y$, find $\dfrac{d^2y}{dx^2}$ in terms of $x$ and $y$.\soln Differentiate: $3-\dfrac{dy}{dx}$, then substitute.\result{$3-3x+y$.}
\end{example}

\newpage
\sech{Practice}

\level{Core}
\begin{enumerate}[leftmargin=*,itemsep=2pt]
  \item Find the equilibrium solutions of $\dfrac{dy}{dt}=y^2-16$.\work{2.3cm}
  \item A solution of $\dfrac{dy}{dt}=0.4(20-y)$ has $y(0)=5$. Find $\displaystyle\lim_{t\to\infty}y(t)$.\work{2.3cm}
  \item For the same equation, describe the solution with $y(0)=30$.\work{2.3cm}
  \item For $\dfrac{dy}{dx}=x^2+y$, find $\dfrac{d^2y}{dx^2}$ in terms of $x$ and $y$.\work{2.3cm}
\end{enumerate}

\level{Exam level}
\begin{enumerate}[leftmargin=*,itemsep=2pt,start=5]
  \item For $\dfrac{dy}{dt}=\dfrac{y(10-y)}{5}$ with $0<y<10$, at what value of $y$ is the solution growing fastest, and how fast?\work{2.8cm}
  \item For the equation in Question 5, at what value of $y$ does a solution have a point of inflection? Give a reason.\work{2.8cm}
  \item A solution of $\dfrac{dy}{dx}=x-y^2$ passes through $(4,2)$. Show that it has a relative minimum there.\work{2.8cm}
  \item For $\dfrac{dy}{dx}=3x-y$, describe the region of the plane in which solution curves are concave up.\work{2.8cm}
\end{enumerate}

\begin{frbox}{Question 9 \textperiodcentered{} Free-response practice}
A fish population $P$, in thousands, satisfies $\dfrac{dP}{dt}=0.2P\left(1-\dfrac{P}{50}\right)$, where $t$ is in years.
\begin{enumerate}[label=(\alph*),leftmargin=*,itemsep=2pt]
  \item Find the equilibrium solutions.\work{2.6cm}
  \item If $P(0)=10$, find $\displaystyle\lim_{t\to\infty}P(t)$ and explain why $P$ is increasing for all $t$.\work{2.6cm}
  \item Find the population at which the fish are increasing fastest, and the rate of increase then.\work{2.6cm}
\end{enumerate}
\end{frbox}

\newpage
\begin{solbox}
\begin{enumerate}[leftmargin=*,itemsep=4pt]
  \item $y=4$ and $y=-4$.
  \item $20$.
  \item It decreases, approaching 20 from above.
  \item $2x+\dfrac{dy}{dx}=2x+x^2+y$.
  \item At $y=5$, halfway between the equilibria; the rate is $\dfrac{5(5)}{5}=5$.
  \item At $y=5$: $\dfrac{d^2y}{dt^2}=\dfrac{10-2y}{5}\cdot\dfrac{dy}{dt}$ changes sign there.
  \item $\dfrac{dy}{dx}=4-4=0$, and $\dfrac{d^2y}{dx^2}=1-2y\dfrac{dy}{dx}=1>0$.
  \item $\dfrac{d^2y}{dx^2}=3-3x+y>0$: above the line $y=3x-3$.
  \item \begin{enumerate}[label=(\alph*),leftmargin=*,itemsep=3pt]
    \item $P=0$ and $P=50$.
    \item $50$. For $0<P<50$ both factors are positive, so $\dfrac{dP}{dt}>0$, and the solution cannot cross the equilibrium $P=50$.
    \item At $P=25$ thousand: $0.2(25)(0.5)=2.5$ thousand fish per year.
  \end{enumerate}
\end{enumerate}
\end{solbox}

\end{document}
